\section{Software utilities for \ROS: dua-utils}
\label{sec:dua-utils}
With the system abstraction layer and the unit template finalized, a single problem is left to solve. It is of practical nature, and arises only when one starts to use the framework and, in particular, the middleware tools integrated into it for robotic software development, such as \ROS (see \Cref{sec:ros2}). The features of the \ROS middleware compose a set of services that the robotic software developer can use to build applications and modules of various degrees of complexity, and the open-source nature of the framework has allowed for a wide community to improve it over the years, reviewing existing functionalities and adding new ones. Nonetheless, the \ROS middleware is a complex system, and its many features can be overwhelming for a beginner or a developer who is not familiar with all of its functionalities. Moreover, some of its APIs are not as intuitive as one would expect, requiring complex code to be written to perform simple tasks which cover instead the vast majority of their use cases.

The \ROS community is aware of these issues, and the development effort is improving the quality of the software every day. However, over some years of experience and active development of \ROS modules, the necessity arose to create a set of utilities that would simplify some parts of the framework, wrapping its complexity and providing a more intuitive interface to the developer. While waiting to propose the same set of changes to the community and have them integrated into the official \ROS distribution, a process that can notoriously happen but also take a long time, the author preferred to develop the same set of utilities as packages that could be integrated in the system configuration provided by every DUA base unit, so that they could be immediately distributed and used in the development of robotic software modules based on DUA. Note that these software modules, as shown later on, are maintained as independent projects with no specific ties to DUA, so that every published work that uses them can be easily integrated into other systems, frameworks, or architectures for the benefit of the scientific community; this is also the case of the two example modules presented later on in this chapter.

This led to the creation of the third foundational pillar of DUA: the software collection called \texttt{dua-utils} \cite{masocco_dua-utils_nodate}. Note that, although at the time of writing it is a collection of utilities for \ROS and robotic software development, it may host different kinds of third-party software packages that can be useful in the development of distributed systems, as long as it is open-source and can be built with colcon. This can be seen as a limitation but it is only due to how it is currently integrated into DUA base units: at the very end of their setup process, the last stage clones the repository into the image filesystem and builds the packages contained in it, cleaning up build artifacts afterwards. This process has two main advantages: first, when the modules are built in this way they end up being part of the DUA base unit image, so that they can be used in the development of software modules based on DUA without the need to install them manually; secondly, the process is automated but performed at the end of every base unit, so if one of the utilities in \texttt{dua-utils} is updated the base unit image can be rebuilt by performing only this last step, reusing previous cached layers and saving time and bandwidth. What described so far may of course change in the future, should the necessity arise to include in \texttt{dua-utils} software packages that cannot be built with colcon or that simply concern other domains of distributed systems development. The motivating idea is that \texttt{dua-utils} is a way to customize the very base layer of the DUA framework which is provided to and intended for its users, who can actively suggest and modify its contents to adapt the framework to their most frequent necessities and usage scenarios.

To give an idea of the set of improvements that DUA provides to the \ROS middleware, as well as a comprehensive overview of the work that has been done so far not only with DUA but also on it, the following sections introduce some very representative utilities that are part of the \texttt{dua-utils} collection.

\subsection{Managing \ROS processes: dua\_app\_management}
\label{sec:dua_app_management}
\texttt{dua\_app\_management} is a collection of software modules and source files for the management of applications and processes in the Distributed Unified Architecture. This package contains modules, both compiled and source-only, to be included in DUA-based projects. It is particularly aimed at the development of \ROS applications, with the aim of simplifying some common tasks and wrapping boilerplate\footnote{In the context of software development, \emph{boilerplate} code refers to long code portions that must be written each time in the same way, to perform what is usually a simple and frequent operation. Programmers usually prefer either simpler APIs or wrappers, \emph{i.e.}, equivalent code items that take less options but are simpler to write, addressing the most frequent use cases.} code, especially that which involves the lifecycle of a process and requires calling the OS to do so.

The package currently contains the \texttt{ros2\_app\_manager} C++ package, which provides two libraries and the implementation of a component container (see \Cref{sec:ros2_nodes}). The first goes by the name of the package itself, and is a header-only library that implements a wrapper object for the usual main thread code of \ROS applications, \emph{i.e.}, processes that run single nodes or collections thereof or act as component containers. The necessity for this library comes from the fact that, if one wants to write the \texttt{main} function for a \ROS application going a little beyond what the examples from the official documentation show, one finds after some time that the same code is written over and over again, with only the node name and the node class changing. Operations performed by this code usually include initializing the ROS MiddleWare context, which is the entrypoint towards the underlying communication middleware used by \ROS, configuring options for the \ROS node, creating a proper executor for the application, and finally adding nodes to it and spinning it. The lifecycle of such an application ends when an asynchronous event, such a signal, occurs and it must shut down, stopping the executor and destroying all such objects in reverse order to correctly release resources and terminate middleware instances. The \texttt{ros2\_app\_manager} library provides a simple API to perform all these operations:
\begin{itemize}
  \item The constructor of the \texttt{ROS2AppManager} class does all the aforementioned initializations, as well as disabling I/O streams buffering for the current process to reduce latencies, which is a common practice in \ROS codes.
  \item The \texttt{ROS2AppManager::run} method spins the executor, and returns only when a signal is caught, which is the usual way to handle the termination of a \ROS application.
  \item The \texttt{ROS2AppManager::shutdown} method stops the executor and releases all resources as previously described.
\end{itemize}
The class is implemented as a template class whose template parameter is the \ROS executor class to use. The library also includes the relevant system headers.

Since signals are the most common way, at least on Linux systems which are the focus of DUA, to request the termination of a \ROS process, the user would benefit from a finer-grained level of control over their action. To clarify the description of what follows, let us briefly introduce the concept of POSIX signals. In operating systems design and programming, and in particular for systems that follow the POSIX standard, \emph{signals} are a way to notify asynchronous events to userspace\footnote{In this context, the term \emph{userspace} denotes application programs that are run and managed by an operating system kernel.} programs, a sort of software interrupt that can be delivered to a process by the kernel either on its own behalf or on behalf of another process, which may even do so due to the user's request. Different signals are identified by numerical codes, and represent different situations and errors. For example, when we press Ctrl-C to stop a program running in a terminal, the terminal tracks that key combination and sends the \texttt{SIGINT} signal to the process it is running, which is the default signal for interrupting a process. The process can then catch the signal and perform some cleanup operations before terminating, provided that it includes a handler for that signal in its code.

In \ROS applications, signals are not extensively used for two reasons: first, applications running on a robot should usually run as long as the robot is on, and second, they require a signal handler to be installed in the application code and some knowledge of the OS APIs. For this reason, \ROS usually handles signals by itself, which is good and handy but limiting in some cases. A first example scenario is if debugging is being performed to track a bug that crashes the applications: appropriately coded signal handlers have access to information concerning what was the memory state when the crash occurred, what kind of error occurred, what was the state of the application and of buses it was accessing, and so on. A second scenario is that of a prototype robot which is being tested, a situation in which applications are usually started and stopped via remote shells opened by users. Suppose that the robot is, \emph{e.g.}, a flying drone, and that such applications allow it to fly safely: it would not be desirable if, \emph{e.g.}, the WiFi link crashes and the drone follows as well because the shells that the applications were running within close abruptly; in this case, ignoring some signals like \texttt{SIGHUP} or \texttt{SIGPIPE} would be a good idea.

The second library of this package, called \texttt{ros2\_signal\_handler}, provides a POSIX-compliant signal handler module for \ROS applications, which replaces the default signal handler installed by the \ROS middleware. The library is a header-only one, and provides a class called \texttt{SignalHandler} implemented as a singleton\footnote{In object-oriented programming, the \emph{singleton} pattern is a software design pattern \cite{freeman_head_2020} that restricts the instantiation of a class to a singular instance within an entire application. Such instance is usually \emph{static}, \emph{i.e.}, it is created, and thus not constructed, by the compiler within the data segment of an executable and is already present in memory when the execution of the program begins.}. It uses the \texttt{sigaction} API to manage handlers, allowing fine-grained control over the signals that are caught and the actions that are performed when they are caught, and can be given access to the executor running in the application to stop it when a signal is caught. The library also includes the relevant system headers. The following methods are provided:
\begin{itemize}
  \item \texttt{SignalHandler::get\_global\_signal\_handler} returns a reference to the singleton instance of the class.
  \item \texttt{SignalHandler::init} installs the signal handler, which is a static method that must be called before any other method of the class. This handler can call user-defined functions when a signal is caught, and can be configured to ignore some signals.
  \item \texttt{SignalHandler::install} to install a handler for a specific signal.
  \item \texttt{SignalHandler::ignore} to ignore a specific signal.
  \item \texttt{SignalHandler::uninstall} to restore the default handler for a specific signal.
  \item \texttt{SignalHandler::fini} to reset the \texttt{SignalHandler} instance to its initial state and restore the default signal handlers for all signals.
\end{itemize}

All these features, although useful, would be useful only when a \ROS application is run as a standalone process, which is not the case for all robotics applications. In reality, what happens is that nodes may also be loaded into a process at run time: such process acts as a component container and nodes are denoted as components, as explained in \Cref{sec:ros2_nodes}. In this case, the container application is usually already compiled and provided as part of the \ROS installation. This is where the third part of this package comes in: the \texttt{dua\_component\_container\_mt} application. It is essentially a component container based on the \texttt{MultiThreadedExecutor} executor class, just like the standard \texttt{component\_container\_mt} provided by \ROS, but with the addition of the \texttt{ros2\_app\_manager} and \texttt{ros2\_signal\_handler} libraries. The application is intended to be used while testing \ROS components which may be dynamically loaded and unloaded at run time on real robot prototypes, thus accounting for the potential issues such as the one described in the example scenario above. For this reason, this container application uses the \texttt{SignalHandler} object to handle many signals and ignoring some of them, and the \texttt{ROS2AppManager} object to manage the lifecycle of the application.

\subsection{Handling Quality of Service in the right way: dua\_qos}
\label{sec:dua_qos}
In \Cref{sec:dds,sec:ros2}, the concept of Quality of Service (QoS) has been presented as a way to configure the behavior of the communication middleware in a distributed system. In \ROS, QoS policies are used to configure the behavior of the underlying communication middleware, being it either the DDS or Zenoh (see \Cref{sec:zenoh}) or something else, and are applied to topics to specify the reliability, durability, and other properties of the communication. The \ROS middleware provides a set of predefined QoS profiles that can be used to configure the communication, but it also allows the user to define custom QoS profiles to meet specific requirements. This system is very flexible and powerful, but often used in an inconsistent way. First, note that \ROS libraries provide default QoS profiles whose names range from \texttt{qos\_profile\_default} to \texttt{sensor\_data}, also explaining those purposes. However, these profiles as well as the rationales behind them come from the ROS 1 age, in which the communication was implemented by the middleware itself on top of UDP and TCP. Since there were mostly only these two protocols to choose from, one had almost always to distinguish between a reliable and an unreliable communication, which led to the establishment of a set of rules of thumb like, \emph{e.g.}, the common belief that data produced continuously by a sensor should be sent with a best-effort policy, since packet loss is always allowed. This is the case unless algorithms processing that data may also account for late packets, in which case a reliable policy should be preferred. Moreover, with the plethora of options that DDS QoS offers, these binary distinctions do not apply anymore. The \ROS API is quite flexible and allows the user to define custom QoS profiles, but the user is left alone to decide which profile to use and how to configure it. This can lead to inconsistent configurations across different nodes and topics, which can result in unexpected behavior and performance issues.

More configuration options imply that a finer control over the transmission policies can be achieved, accounting for an issue that ROS 1 struggled with: intrinsic difference in data types being transferred. For example, an inertial sensor may produce many samples every second but they will all correspond to small packets since they are made of a few numeric fields. Conversely, a LiDAR scan may have a lower frequency but produce large packets, since each scan is made of many points; a similar situation may arise with cameras. Reliability is not the only issue to consider: the frequency of the data, the size of the packets, the importance of the data, and the network conditions are all factors that can influence the choice of the QoS profile to use.

The \texttt{dua\_qos} library \cite{masocco_dua_qos_nodate} provides API that return QoS profiles that account for these issues, distinguishing first between the desired reliability of the communication and then between the data type being transferred. The repository contains two software libraries: \texttt{dua\_qos\_cpp} for the C++ language and \texttt{dua\_qos\_py} for the Python language. Apart from their actual implementation, whose details concern only the syntax of each language, the functionalities they offer are the same. Before going further, note that DUA does not enforce the use of these QoS profiles in any way nor place: they are only included as a suggestion for the user to simplify the development of \ROS applications, and are also designed to be compatible as much as possible with the QoS policies that are both provided by the \ROS middleware and used in the majority of the \ROS software available today. Lastly, the profiles are intended to be used only with topics. One can modify QoS for internal communications of services and actions too, but that is usually not recommended: those primitives are synchronous and reliable by nature and, in the case of actions, require specific QoS policies to achieve stateful communication. For these reasons, the \texttt{dua\_qos} libraries provide only QoS profiles for topics.

The first two categories are \texttt{Reliable} and \texttt{BestEffort}: these map to two distinct libraries in both packages, whose contents are almost the same. The difference they enforce is the \texttt{ReliabilityPolicyKind}, \emph{i.e.}, the reliability part of the QoS policies provided by each library. There is no additional distinction between the two categories since it is actually up to the user to decide which one to use: today, the \texttt{Reliable} category applies to practically every transmission since network congestion and bandwidth optimization are managed independently by the communication middleware used by \ROS, mainly thanks to the \texttt{HistoryPolicy}. This policy is set to \texttt{KEEP\_LAST} by default, which means that only a fixed-size buffer of samples is kept in memory to be either sent or received, and that the rest is discarded, independently of the configured reliability which comes into play shortly after, when samples are actually being transmitted; this mechanism, although not sophisticated, acts as a sort of congestion control that works well in most cases. All APIs in the \texttt{dua\_qos} libraries allow to set the depth of such buffers, but provide a default value that is usually good enough for most cases.

For each package, both \texttt{Reliable} and \texttt{BestEffort} libraries then provide QoS profiles via the following APIs:
\begin{itemize}
  \item \texttt{get\_datum\_qos}: returns a QoS profile for generic data, \emph{i.e.}, data that is not particularly large in size or requires special handling on the QoS size. This is due to the above reasoning, thanks to which both commands and setpoints, and sensor data are, \emph{e.g.}, the same thing.
  \item \texttt{get\_scan\_qos}: returns a QoS profile for LiDAR scans, spatial scans, or similar data. The peculiarity of this data is that while it may have relatively high or low frequencies, it gets easily very large in size, so a low history depth is suggested here to minimize the risk of network congestion and memory exhaustion.
  \item \texttt{get\_image\_qos}: returns a QoS profile for images, whose size may vary from small to very large, but whose frequency rarely exceeds a few tens of Hertz. The history depth default value is not as low as for scans, but still lower than for generic data and configurable by the user.
\end{itemize}
The only difference among the two libraries is that the \texttt{Reliable} one also provides the \texttt{get\_persistent\_qos}, a particular type of policy intended for reliable transmissions that may be resumed at a later point in time. This is possible thanks to a feature of the DDS standard called the \texttt{DurabilityPolicy}, discussed in \Cref{sec:dds}, which allows a publisher to store all samples in a persistent storage and send them again them when a new subscriber appears. This is useful for data that may not be produced continuously but is still important, like configuration data or state information, and is used in this library together with a higher history depth to ensure that all samples are stored and sent.

\subsection{Improving node development: params\_manager, dua\_node}
\label{sec:params_manager_dua_node}
Among the many \ROS features (see \Cref{sec:ros2}), there is one of paramount important concerning the nodes: the possibility of specifying configurable variables for them that can be inspected and set at any point in time, known as \emph{node parameters}. These parameters are usually set at launch time, either by the user or by the launch system via configuration files, and are then read by the node to configure its behavior. The \ROS middleware provides a set of APIs to manage these parameters, whose task is that of building and storing in each node a database of the parameters that it has. The keys used in this database are strings, which contain the name of a parameter. Each parameter is then associated with its type, which can be a string, an integer, a floating-point number, a boolean, or a list of any of these types, and its current value. There are also other flags and fields, identifying whether a parameter is read-only or, in case it is of numeric type, its valid ranges.

\ROS configuration files allow the user to specify only parameter values, meaning that all the aforementioned information must be provided when the parameter is \emph{declared}. In \ROS code, a \emph{parameter declaration} is the act of calling an API that creates an entry into the parameter database of the node, writing all such information which is passed as arguments to the API. To correctly declare parameters passing all the correct and relevant information, one has to write quite the quantity of code. This is not a problem when the number of parameters is small, but it can become cumbersome when the number of parameters grows, especially if they are of different types and have different constraints. As mentioned earlier, this is a classic example of boilerplate code.

Another issue concerns parameter updates and their usage in the actual code of the node. Each time a parameter must be read, the node must call an API that retrieves the parameter from the database, and then cast it to the correct type. This is not a great concern if the parameter is read only once or with a low frequency but when it is, \emph{e.g.}, the gain of a control law that must be evaluated at each high-frequency iteration of a control loop, the overhead of the database queries might become significant. Assuming that the parameter type is never going to change, something that the \ROS APIs allow but is seldom necessary in practice, it would be preferable to store the parameter value in the node and update it only when the parameter is updated. This way, parameter reads would be much faster. To achieve this, a callback must be registered with the \ROS node that is called when an update is requested for a set of parameters: when this happens, some checks like the range of the parameter have already been performed but other ones like the type of the parameter have not. Moreover, there can be only one such callback in a node, which means that it must check for which parameter an update has been requested and act accordingly.

The amount of code that one has to write to achieve the aforementioned functionalities, in every \ROS software and every time in the same way, grows uncontrollably with respect to the amount of parameters that each node has. The code that manages parameters in a \ROS node becomes unnecessarily long and complex, difficult to manage and debug. This is the motivation behind the development of the \texttt{params\_manager} package \cite{masocco_params_manager_nodate}, which provides a set of utilities to simplify the management of parameters in \ROS nodes. This package currently contains only an implementation for C++.

Before the \texttt{params\_manager} C++ library comes into play, there is a Python script that can be used to automatically generate the code that declares the parameters in the \ROS node class. This script is automatically called by a CMake directive that extends the ament installation (see \Cref{sec:ros2_buildsys}). Thus, the first step, as shown in the documentation of \cite{masocco_params_manager_nodate}, is to add a YAML file to the source code of the node that, with a specific syntax shown in \Cref{lst:paramsmanager}, contains the list of parameters that the node must manage together with all their metadata. The script will then parse this file and generate a C++ source file that gets added to the build process, containing a routine called \texttt{init\_parameters} that declares all parameters in the node and sets up the update callback. To make this work, a bit of code must be added to the node class, as shown in the documentation: a \texttt{Manager} object must be instantiated as a member of the node class, and the \texttt{init\_parameters} function must be declared as a method of the node and called in its constructor, better if at its beginning since from that point on all parameters will have been declared, set, and ready for usage. This way, the boilerplate code is reduced to a bare minimum of five lines, and all relevant parameter data must be written in a configuration file which is much easier to manage and debug than the code itself.
\begin{center}
\begin{minipage}[t]{.9\textwidth}
\begin{lstlisting}[caption={Example of \texttt{params\_manager} YAML configuration file \cite{masocco_params_manager_nodate}.}, label={lst:paramsmanager}]
header_include_path: dir/node_header.hpp
namespace: node_namespace # This is optional
manager_name: manager # This is optional
node_class_name: MyNode

params:
  camera_offset:
    type: double
    default_value: 0.1
    min_value: 0.0
    max_value: 0.5
    step: 0.0
    description: "Distance between camera and drone center."
    constraints: "Must be in meters."
    read_only: false
    validator: function_name3 # This is optional

  hsv_from:
    type: integer_array
    default_value:
      - 0
      - 1
    min_value: 0
    max_value: 255
    step: 1
    description: "HSV color lower range."
    constraints: "."
    read_only: false
    var_name: hsv_from_ # This is optional
\end{lstlisting}
\end{minipage}
\end{center}

\Cref{lst:paramsmanager} shows also two other useful features of the \texttt{params\_manager} package. The first is the possibility to specify a \emph{validator} routine or a variable name: the autogenerated parameter update callback checks the name and type of a parameter and, if a match is found, may perform other operations before approving the update. One of these is to perform an additional check, whose outcome decides if the parameter update is approved, which may be coded by the user in a function whose name is specified in the YAML file, informally referred to as a \emph{validator}. Alternatively, to account for the necessity to quickly store the new value directly inside the node as explained above, the user can specify the name of a variable member of the node class that the callback will access and overwrite with the new parameter value. Of course, the two things can be done at the same time: it is up to the user in the first case to code the variable member update inside the validator routine. To keep track of each parameter's metadata together with these two additional values, the \texttt{params\_manager} library uses the \texttt{std::set} data structure from the \texttt{<set>} C++ standard library header, which is implemented as a red-black tree for maximum efficiency \cite{guibas_dichromatic_1978}.

After the \texttt{params\_manager} library was implemented and successfully tested, the question arose whether it would be possible to eliminate the necessity of the user to write the five lines of code that instantiate the \texttt{Manager} object, call the \texttt{init\_parameters} function in the node constructor, and then clear the \texttt{Manager} at the end. Hiding the management of the \texttt{Manager} would be possible in case the base \ROS node class that each node is derived from already included an object of the \texttt{Manager} class as a variable member. In the \texttt{dua\_node} package \cite{masocco_dua_node_nodate}, an alternative \ROS base node class is provided that includes the \texttt{Manager} object and handles its lifecycle together with that of the node. This way, the only necessary lines of code to manage parameters in a \ROS node are those that declare and call the \texttt{init\_parameters} function; this is due to the fact that, by the current implementation, \texttt{init\_parameters} is a member function of the derived class and can only be called by its constructor, after the base class one has run. The \texttt{dua\_node} package is intended to be used in conjunction with the \texttt{params\_manager} package, and is included in the \texttt{dua-utils} collection.

\subsection{Wrapping clients: simple\_serviceclient, simple\_actionclient}
\label{sec:simple_serviceclient_actionclient}
\Cref{sec:ros2} presented ROS, in both its first and second version, as a project aimed at providing a middleware to robotic software developers that offers many services, the most important of which is distributed communication. To this end, \ROS provides three communication primitives: topics, services, and actions. The last two are synchronous communication primitives, meant to implement client-server communication over a short time span in the former case, and a possibly longer one in the latter case, with the additional feature of being stateful. The \ROS middleware provides APIs to create and manage services and actions, which are grouped in the server and client sides. Servers are usually simple to implement: the user instantiates the server object and defines the relevant callbacks for it. In the case of actions there are usually multiple callbacks involved, to account for every possible asynchronous event requested by the client or terminal state (see \Cref{fig:ros2act,fig:ros2goalfsm}), but the developer can choose which ones to implement based on which features of the communication their code shall actually support.

For both services and actions, the client side is instead a completely different beast. In both C++ and Python, client-side APIs are based on asynchronous programming, which is a programming paradigm that allows the user to perform multiple tasks concurrently, without the need to wait for the completion of one before starting the next. This is achieved by using callbacks together with \emph{futures}, which can be essentially described as variables whose value becomes available at a later point in time. Things get increasingly complicated as one steps from service clients to action clients, so this is the order in which the two situations are addressed here.

The developer writing code making use of a service client must first instantiate such object, then send a service request through it. The result of this API call is a future, which will eventually hold the value of the service response. The developer can act in three ways here, all made possible by the \ROS API: they can block and synchronously wait for the future to get a value, they can attach a callback to the future that will be called when the value is available, or they can poll the future at any time they want in the code and check for its value. The API to manage each of these cases is not particularly complex, but it can get quite verbose. As for all the previous sections describing packages of \texttt{dua\_utils}, the next one is motivated by a necessity that arises when one writes actual \ROS software: in the vast majority of cases, if you call a service, you need the response data right after, and you need it to be available before the next line of code is executed. This is why the first eventuality among the three described above is the most common, but it still requires the developer to write a lot of code to manage the future, put the node in a synchronous wait state by spinning an executor, and then get the value from the future. This is the reason why the \texttt{simple\_serviceclient} package \cite{masocco_simple_serviceclient_nodate} was developed.

This package offers two equivalent libraries for both C++ and Python, which offer the same methods with the exact same signatures in both languages. The libraries implement a client class that wraps the original \ROS service client, extending its functionalities. Since the other situations are better handled with the official \ROS APIs, the \texttt{simple\_serviceclient} libraries offer only two methods:
\begin{itemize}
  \item \texttt{call\_sync}: this method sends a service request and waits for the response, returning it when it is available. The method is blocking, and takes a flag argument telling whether the node must be added to an executor or not; this is to support the situation in which the node is already added to an executor and callbacks execute in separate threads, so waiting for the future consists in blocking the calling thread waiting for it to get a value using the future object API. In both scenarios, the only code that can execute while the future is being waited for is node callbacks.
  \item \texttt{call\_async}: this method sends a service request and returns a future that will hold the response when it is available. The method is non-blocking, and provided for convenience.
\end{itemize}
In C++, these libraries contain template classes whose template parameter is the type of service, so that the code can automatically handle the service request and response types.

The wider variety of cases that the stateful communication provided by actions defines is the main reason behind the complexity of its client-side. Essentially, each operation that the client requests to the server is implemented as a service call, and must be treated in the code in a way similar to what has been described before, \emph{i.e.}, using futures and eventually callbacks. Note that even the most desirable outcome for an action from the perspective of the client, \emph{i.e.}, completing a goal, requires two service calls: one to request the acceptance of the goal and one to request and get back the result (see, again, \Cref{fig:ros2goalfsm}). This easily leads again not only to a lot of boilerplate code, but also to complex code that is difficult to manage and debug. Moreover, there is an additional problem: the \ROS action API is not as well organized as the service API, due to its history. It has been one of the last features to be ported, regarding the communication services, from ROS 1 to \ROS, and its implementation is as such split among two sets of libraries: the language client library, that is \texttt{rclcpp} for C++ and \texttt{rclpy} for Python, and the \texttt{rclcpp\_action} and \texttt{rclpy.action} libraries. This dichotomy, still unresolved, has led to a situation in which the style of the client API might differ even significantly between its different parts, which reduces clarity and increases the learning curve for the developer. One of the most evident examples of this is the numerous types that are involved in client-side action programming, which model goals, their handles, futures, feedbacks, etc., and their signatures do not follow a consistent style. The API is very powerful and allows one to achieve very complex asynchronous, event-based behaviors to handle scenarios of various degrees of complexity, but the problem is the same as before. In reality, there are mostly two ways in which actions are called: either the client sends a goal and waits immediately after for the result, or the client sends a goal and then either polls for the result or requests a goal cancellation. The \texttt{simple\_actionclient} package \cite{masocco_simple_actionclient_nodate} was developed to address these issues, providing an action client class that acts as a wrapper for the original action client. The package contains again two libraries, one for C++ and one for Python, whose contents and implementations are the same. For simplicity, the following explanation will refer only to the C++ library, which implements a template class.

The first big improvement of this class over the original API is that it completely hides all the type definitions that are required to handle an action client: the only template parameter is the action type. This leads to a very simple instantiation of an action client, where all type definitions are handled internally by the class. The class constructor can also take feedback callback as argument, to be called when feedback messages are received from the server. Then, the class implements three sets of methods, whose names are self-explanatory and that implement the most important moments of an action call lifecycle. For each, this library also provides both synchronous and asynchronous variants, which counts to six methods in total:
\begin{itemize}
  \item \texttt{send\_goal}: sends a goal to the server, and returns a future that will hold the result of the goal request.
  \item \texttt{send\_goal\_sync}: sends a goal to the server, and returns the result of the goal request.
  \item \texttt{cancel}: cancels the goal that is passed as argument and returns a future that will hold the result of the goal cancellation.
  \item \texttt{cancel\_sync}: cancels the goal that is passed as argument, and returns the result of the goal cancellation.
  \item \texttt{get\_result}: requests the result of the goal that is passed as argument, and returns a future that will hold the result of the goal.
  \item \texttt{get\_result\_sync}: requests the result of the goal that is passed as argument, and returns the result of the goal.
\end{itemize}

Finally, the \texttt{simple\_actionclient} library also provides a method to handle the entire lifecycle of an action call in one go, from the goal request to the reception of the result: the \texttt{call\_sync} method. The signature of this method for the C++ library is shown in \Cref{lst:callsync}. The first thing to note about it are its arguments: apart from the goal to send, it takes a flag and a series of timeouts. The flag tells the method whether to cancel the goal if a timeout occurs, and the timeouts are the maximum time that the method will wait for the goal request to be sent, for the result to be received, and for the goal to be cancelled. The method returns a tuple, implemented in C++ with the \texttt{std::tuple} standard library class, with three elements:
\begin{itemize}
  \item a boolean flag that tells if the goal has been accepted or not by the server;
  \item a \texttt{ResultCode} object, which is an enumeration defined in the \texttt{rclcpp\_action} library that tells the result of the goal request, \emph{e.g.}, if the goal has been accepted, rejected, or cancelled;
  \item a shared pointer to the result of the goal, which can be valid or not depending on whether it was returned by the server, and that has type \texttt{ActionResultT} which is derived from the action type, the template parameter of the class.
\end{itemize}
Internally, the method performs each step by requesting each operation to the server and then synchronously waiting for the response, spinning the node. As the call progresses, the method checks for timeouts and cancels the goal if necessary, handling all possible outcomes and return values. When the return value of this method is accessed, by the way it is designed as a tuple holding the three values discussed above, the caller can immediately determine how the action call went since those three values together encode every possible outcome of the sequence of operations that the method performs. This method is the most powerful and flexible of the library, and is intended to be used in the vast majority of cases where an action is called. The other methods are provided for convenience, and to allow the user to handle each step of the action call lifecycle separately if needed. For more advanced use cases, the original \ROS API is again the best choice.
\begin{center}
\begin{minipage}[t]{.9\textwidth}
\begin{lstlisting}[caption={Signature of the \texttt{Client::call\_sync} class method from the C++ action client wrapper library \texttt{simple\_actionclient\_cpp} \cite{masocco_simple_actionclient_nodate}.}, label={lst:callsync}]
std::tuple<bool,rclcpp_action::ResultCode,std::shared_ptr<ActionResultT>>
  call_sync(
    const typename ActionT::Goal & goal_msg,
    bool cancel_on_timeout = false,
    int64_t send_goal_timeout_msec = 0,
    int64_t get_result_timeout_msec = 0,
    int64_t cancel_timeout_msec = 0);
\end{lstlisting}
\end{minipage}
\end{center}

\subsection{Implementing \ROS automata: transitions\_ros}
\label{sec:transitions_ros}
As illustrated in \Cref{sec:prb_design_challenges}, the software that operates a complex robot typically follows a hierarchical organization. The lower levels usually consist of firmware, controllers, and device drivers. More advanced modules rely on the abstractions provided by these lower layers to perform sophisticated operations. At the top, a single module typically acts as a coordinator, ensuring the proper functioning of the entire system and offering interfaces to human operators. This module usually requires complete access to the software stack of the robot for control purposes, and its design must be sophisticated enough to handle the diverse situations that might arise during the execution of the prescribed tasks that the robot must perform. The higher the level of autonomy required, the more complex this kind of software is going to be.

Finite-state machines (FSMs) are a renowned tool in computer science \cite{wagner_modeling_2006}: their mathematical foundation derived from graph theory describes a system whose behavior is represented by a sequence of discrete configurations within a finite set, called \emph{states}. These states are interconnected by a collection of \emph{transitions}, which may lead to another state or loop over the same one. Transitions occur when specific conditions are met, which alter the current state of the machine. FSMs have been widely utilized in modeling various hardware and software systems, ranging from basic algorithms and communication protocols \cite{brand_communicating_1983} to digital control systems \cite{sklyarov_hierarchical_1999, pola_symbolic_2016} and distributed systems \cite{schneider_implementing_1990}.

Before moving forward, it is essential to clarify the type of FSM relevant to this use case. In software engineering, two main categories of FSMs are recognized. The first category, often referred to as a \emph{passive} FSM, is the more commonly used of the two. This type mainly serves as an internal representation of the state of an object or algorithm, acting as a data structure that holds information about the current state without performing any actions. Methods that interact with a passive FSM are called externally and their primary role is to keep a record of the execution state of the program based on the information encoded within it. The module presented in this section focuses on the second category of FSM, known as an \emph{active} FSM. Unlike its passive counterpart, an active FSM is an autonomous entity that retains state information and actively executes operations on data and interfaces with other system components by calling appropriate APIs. An active FSM plays a crucial role in autonomous robotics by governing the behavior of the robot. This includes decision-making processes, executing movement commands, and interpreting data from its various subsystems. It is one of the most widely used paradigms for developing supervisory software for autonomous robots, as it provides a structured approach to managing the complexity of the robot's behavior; it is not the only one though, since its limitations start to show up when the automaton's complexity grows and a more modular and intrinsically nested structure would be preferable, such as behavior trees \cite{iovino_survey_2022}.

This discussion clarifies that, when looking forward to develop a supervisory software for an autonomous robot that can be modeled as a finite-state machine, an active FSM is the most appropriate choice. Moreover, if the software running on the robot is based on \ROS, having an active FSM implementation which is designed to directly support \ROS communication primitives and other middleware functionalities would be highly beneficial. The \texttt{transitions\_ros} package \cite{masocco_transitions_ros_2022} was developed to address this need: it provides the implementation of an active FSM capable of seamless interaction with \ROS software, thereby enabling sophisticated autonomous control functionalities. Such an active FSM is intended to facilitate high-level autonomy in robotic systems, allowing them to operate effectively in dynamic environments by integrating sensory inputs, executing control logic, and interfacing with other modules cohesively.

When dealing with a \ROS-based robot, one can broadly isolate three types of tasks that an automaton must handle:
\begin{itemize}
  \item \textbf{Asynchronous:} the robot must react to some external events that do not alter its state but require data processing, \emph{e.g.}, detecting a target.
  \item \textbf{Synchronous:} the robot must execute operations specific to its current state, \emph{e.g.}, arming, takeoff, landing, or exploring the environment.
  \item \textbf{Collaborative:} the robot must perform complex operations requiring real-time coordination with another agent or human operator.
\end{itemize}
Given the three communication primitives offered by \ROS (see \Cref{sec:ros2}), a logical approach for implementing a \ROS-enabled FSM would be to handle asynchronous tasks through appropriate topic subscription callbacks. In contrast, services and actions would be ideal to manage synchronous and collaborative tasks. The FSM implementation should thus include a reference to a \ROS node object, acting as a communication endpoint towards the middleware. To ensure the FSM is active in the sense described above, each state should be associated with a specific \emph{routine}, \emph{i.e.}, a function executed upon entry to the state, whose return value determines the subsequent transition. Furthermore, the FSM object should provide a \texttt{run} method that, when invoked by a running thread, sets the automaton in motion within the context of the calling thread, and returns control only once a terminal state is reached and no further transitions are possible.

The \texttt{transitions} \cite{yarkoni_transitions_2017} Python library is one of the most famous and versatile open-source software libraries that implements a finite-state machine. It has many standard features of software FSM implementations, such as support for hierarchical state machines and guard conditions evaluation. However, the only crucial feature here is the possibility of easily defining states and transition tables, operations that the library allows by specifying a dictionary for the former and a nested list for the latter. Strings identify both states and transitions by name, and the library supports regular expressions to make some definitions quicker such as, \emph{e.g.}, catch-all transitions to an error state. With respect to the other use cases described so far in previous sections, the choice of this library which is based on the Python programming language has been motivated by its high-level nature, which reduces the time and code complexity required to encode and test even sophisticated routines, deeming it appropriate for a supervisor software module whose implementation should be as straightforward as possible.

The proposed \texttt{transitions\_ros} library \cite{masocco_transitions_ros_2022} starts by extending the \texttt{State} object of the \texttt{transitions} library, by deriving its class: each state must hold a reference to a \ROS node object, and a function to call as their state routine by the new \texttt{routine} method of the \texttt{State} class. Such a routine must be defined elsewhere in a software module that uses this library to build a specific automaton, must take a \ROS node as argument to interact with the rest of the software architecture, and must return a string that identifies the next transition that must occur based on the outcome of the routine itself; such string should be empty in the case of a terminal state. An example of transitions and state tables for the \texttt{transitions\_ros} library is given in \Cref{lst:transitions}. Next, the library extends the \texttt{Machine} object in a similar way: leveraging a specific use case of the \texttt{transitions} library, in \texttt{transitions\_ros} a machine is created as an independent object that holds information about states and transitions, as well as a \texttt{run} method that sets things in motion. The \texttt{run} method, when called, executes the routine of the initial state, performs subsequent transitions, and calls new state routines based on their string return value, stopping only when an empty string is returned by some routine which indicates that the automaton execution must stop. Finally, to receive and process the incoming data and requests of the node while the automaton is, \emph{e.g.}, synchronously waiting for some task to finish, the \texttt{Machine} object offers the \texttt{wait\_spinning} method that adds the node to an executor, spins it up to a configurable timeout, and then may perform further processing of new data.
\begin{center}
\begin{minipage}[t]{.9\textwidth}
\begin{lstlisting}[caption={Example of transitions and state tables in the \texttt{transitions\_ros} library \cite{masocco_transitions_ros_2022}.}, label={lst:transitions}]
transitions_table = [
    ['START', 'INIT', 'WAIT_FOR'],
    ['READY', 'WAIT_FOR', 'FIRST_TAKEOFF'],
    ['AIRBORNE', 'FIRST_TAKEOFF', 'EXPLORE'],
    # ...
    ['MISSION_DONE', 'DONE', 'RTH'],
    ['AT_HOME', 'RTH', 'COMPLETED'],
    ['ERROR', '*', 'ERROR']
]

states_table = [
    {'name': 'INIT', 'routine': init_routine},
    {'name': 'RTH', 'routine': rth_routine},
    {'name': 'COMPLETED', 'routine': completed_routine},
    {'name': 'ERROR', 'routine': error_routine}
]
\end{lstlisting}
\end{minipage}
\end{center}

A \ROS Python application using this library would have to, in order, define a \ROS node class with all the required interfaces, specify states and transitions tables, define routines for each state, then create both the node and machine objects and call the \texttt{Machine.run} method in a dedicated thread. The proposed implementation keeps the code concisely organized by design, allowing each of the definitions as mentioned above to be performed in a dedicated source file. To conclude, let us explain how this FSM implementation would handle the classes of tasks mentioned before:
\begin{itemize}
  \item \textbf{Asynchronous:} their triggering or processing can be encoded within topic subscription callbacks of the node, which may be executed during a call to \texttt{wait\_spinning} or in a similar busy-wait context and may modify data that influences the automaton flow.
  \item \textbf{Synchronous:} these can be performed entirely within state routines by calling, \emph{e.g.}, service or action clients of the node that invoke specific operations in the underlying software modules, and inspecting their results.
  \item \textbf{Collaborative:} these usually require a mixture of the two previous approaches, since the task execution might be synchronously coordinated by all parties but some of its conditions might arise asynchronously, \emph{e.g.}, the mutual triggering of the task itself among multiple autonomous robots or the human operator.
\end{itemize}
Finally, note that the design of this active FSM implementation integrates very well with the client wrappers for services and actions presented in \Cref{sec:simple_serviceclient_actionclient}.
